\documentclass[11pt]{article}

\usepackage[margin=1in]{geometry}
\usepackage{amsmath,amssymb,amsthm}
\usepackage{mathtools}
\usepackage{booktabs}
\usepackage{microtype}
\usepackage{hyperref}

\IfFileExists{cleveref.sty}{\usepackage[nameinlink,noabbrev]{cleveref}}{}
\providecommand{\cref}[1]{\ref{#1}}
\providecommand{\Cref}[1]{\ref{#1}}

\title{Polynomial Evaluation in the Boolean-Kernel Basis on Modern Processors}
\author{Ali Mkhida\\Algorizk Labs}
\date{}

\begin{document}
\maketitle

\begin{abstract}
Polynomial evaluation is a basic operation in computer algebra and in
several polynomial-based proof systems.  Its cost is commonly described by
the number of arithmetic operations, although on current processors the
dependency structure of the computation, memory traffic, and modular
reduction also affect performance.  In this article, we compare arbitrary
single-point evaluation from three univariate representations: monomial
coefficients, values on a roots-of-unity Lagrange domain, and a
Boolean-kernel basis.  The Boolean-kernel representation admits a binary
reduction in which all folds at a fixed level are independent.  We study
this reduction over the Goldilocks field on an Apple M1 processor and
compare reduction specialized to the Goldilocks modulus with Montgomery
multiplication.  At \(N=2^{20}\), the optimized Boolean-kernel evaluator
takes \(1.0560\) ms, compared with \(1.4116\) ms for a blocked parallel
monomial implementation in the same benchmark run, a reduction of
\(25.2\%\) in wall time.  Isolated Goldilocks and Montgomery
multiplications are nearly indistinguishable, while the complete kernel
evaluator is consistently faster with the specialized arithmetic backend.
The measurements therefore indicate that the observed gain results from
the interaction between the polynomial representation, its dependency
graph, and the field-arithmetic implementation.
\end{abstract}

Polynomial evaluation is among the elementary operations on which more
complex numerical and algebraic software is built.  Classical evaluation
schemes and their interaction with processor parallelism have been studied
extensively; recent processor-oriented comparisons include
Ewart et al.~\cite{ewart2020}.  Horner's rule gives a compact recurrence for
a dense polynomial in monomial form, while divide-and-conquer schemes
expose more independent operations at the price of additional work.  In
applications that store a polynomial by its values on a structured domain,
interpolation formulas provide another route to an arbitrary evaluation
point; barycentric Lagrange interpolation is a standard example
\cite{berrut2004}.

The operation is also present in transparent proof systems and their
polynomial subroutines.  FRI, for example, operates on Reed--Solomon codewords
and repeatedly reduces polynomial degree by structured folding
\cite{fri2018}, while STARK constructions combine such polynomial machinery
with transparent computational-integrity protocols \cite{stark2018}.  In
these settings, a polynomial may be stored in coefficient form, as
evaluations on a domain, or in another representation chosen to simplify a
later protocol stage.  Consequently, the cost of evaluating at a
verifier-derived point depends not only on the degree but also on the
representation in which the polynomial is available.

Arithmetic counts alone do not determine execution time on a modern
processor.  A short dependency chain can permit several independent field
operations to proceed concurrently, whereas a recurrence with fewer
operations may serialize the execution.  The cost of the field
multiplication is another important part of the comparison.  For the
Goldilocks prime
\[
p=2^{64}-2^{32}+1,
\]
one may either exploit the special form of the modulus or keep field
elements in Montgomery representation.  These two choices have different
instruction sequences and need not lead to the same ranking of polynomial
evaluation methods.

In this work we consider a third representation, indexed by Boolean
vectors,
\[
K_y(X)=\prod_{i=0}^{m-1}\left(X^{2^i}+y_i\right),
\qquad
y\in\{0,1\}^m .
\]
For \(N=2^m\), a polynomial of degree less than \(N\) is written
\[
P(X)=\sum_{y\in\{0,1\}^m}c_yK_y(X).
\]
At a fixed point \(x\), two coefficients whose indices differ only in bit
\(i\) can be replaced by one coefficient using
\[
(a,b)\longmapsto x^{2^i}(a+b)+b.
\]
The resulting computation is a complete binary reduction tree.  It has
linear work, but a dependency structure substantially different from
Horner's rule.

The purpose of this study is to compare these evaluation structures on the
same processor and with the same finite-field implementation.  We first
describe the polynomial representations and the arithmetic work associated
with each evaluation method.  We then examine the Goldilocks and Montgomery
multiplication algorithms, describe the processor and measurement method,
and finally present the benchmark results and their relation to the
dependency graphs of the algorithms.

\section{Polynomial Evaluation}
\label{sec:poly-eval}

Let
\[
P(X)=a_0+a_1X+\cdots+a_nX^n,\qquad a_i\in\mathbb{F},
\]
where \(\mathbb{F}\) is a finite field.  A large number of evaluation
schemes can be used for the computation of \(P(x)\), and their arithmetic
complexity need not predict their performance on an out-of-order processor
\cite{ewart2020}.  In this study we retain the schemes that correspond to
the representations used in the benchmark: the classical Horner method, a
blocked Horner method exposing parallelism, evaluation from values on a
roots-of-unity domain, and the Boolean-kernel representation.

As in classical studies of polynomial evaluation, the first comparison is
made in terms of arithmetic operations.  Table~\ref{tab:arith-structure}
summarizes the main costs.  These counts do not yet include the cost of the
modular reduction used to realize a field multiplication.  This point will
be treated separately because the Goldilocks modulus allows more than one
reasonable multiplication strategy.

\begin{table}[ht]
\centering
\caption{Arithmetic structure of the evaluation methods considered in this
study.  For the last two rows \(N=2^m\).}
\label{tab:arith-structure}
\begin{tabular}{llll}
\toprule
Method & Multiplications & Additions & Main dependency\\
\midrule
Horner &
\(n\) &
\(n\) &
chain of length \(n\)\\

Blocked Horner &
\(n+\theta(B)\) &
\(n\) &
independent blocks\\

Lagrange, batch inversion &
\(O(N)\)+one inversion &
\(O(N)\) &
two linear scans\\

Boolean-kernel fold &
\(N-1\) &
\(2(N-1)\) &
binary tree\\
\bottomrule
\end{tabular}
\end{table}

The number of operations gives useful information, but it is not sufficient
to predict the execution time.  Horner's rule has the smallest and most
regular operation count, but its successive multiply-add steps form a
dependency chain.  The blocked Horner and Boolean-kernel methods expose
independent operations.  The Lagrange method requires more arithmetic and
temporary storage, but part of its work is also parallel.  The relation
between these properties and the measured execution time is one of the
questions considered in this article.

\subsection{Horner Method: Classical}
\label{sec:horner}

The classical Horner method evaluates the monomial representation
recursively.  For an input value \(x\),
\[
P(x)
=
a_0+x\left(
a_1+x\left(
a_2+\cdots+x a_n
\right)\right).
\]
Starting from the highest coefficient, the recurrence is
\[
r_n=a_n,\qquad
r_i=a_i+xr_{i+1},
\qquad i=n-1,\ldots,0,
\]
and \(P(x)=r_0\).

The method uses \(n\) field multiplications and \(n\) field additions.  Its
advantage is therefore clear from the arithmetic point of view.  The
drawback is also immediate: \(r_i\) cannot be computed before
\(r_{i+1}\).  The critical path contains all \(n\) multiplication steps.
For large \(n\), the execution of one polynomial gives little
instruction-level parallelism even when the processor has several
execution units available.

\subsection{Horner Method: Blocked}
\label{sec:blocked-horner}

The dependency chain can be shortened by splitting the coefficients into
blocks.  Let \(B\) be a block size and, for simplicity, suppose that
\(N=n+1\) is divisible by \(B\).  Write
\[
P(X)
=
\sum_{j=0}^{q-1} X^{jB}P_j(X),
\qquad
q=\frac{N}{B},
\]
where
\[
P_j(X)
=
a_{jB}+a_{jB+1}X+\cdots+a_{jB+B-1}X^{B-1}.
\]

The \(q\) polynomials \(P_j(x)\) are independent and can be evaluated in
parallel by the classical Horner method.  Once they have been computed,
the final combination is another Horner recurrence, now in
\[
x^B:
\qquad
P(x)
=
P_0(x)
+
x^B\left(
P_1(x)+
x^B\left(
\cdots+
x^B P_{q-1}(x)
\right)\right).
\]

Ignoring the cost \(\theta(B)\) required to obtain \(x^B\), the block
evaluations and the final combination together use \(n\) field
multiplications and \(n\) field additions.  The arithmetic count is thus
close to the classical method, but the critical path is reduced because
the block evaluations can proceed concurrently.  In the implementation
used later, the blocks are distributed between Rayon workers and the final
combination is performed after the local Horner evaluations.

\subsection{Lagrange Representation on a Roots-of-Unity Domain}
\label{sec:lagrange}

A polynomial may also be represented by its values on a structured
evaluation domain.  Let
\[
H=
\{1,\omega,\omega^2,\ldots,\omega^{N-1}\},
\]
where \(\omega\) is a primitive \(N\)-th root of unity, and let
\[
v_i=P(\omega^i).
\]
For a point \(x\notin H\), the Lagrange formula reads
\[
P(x)=
\sum_{i=0}^{N-1}
v_i L_i(x),
\]
with
\[
L_i(x)
=
\frac{(x^N-1)\omega^i}
     {N(x-\omega^i)}.
\]

The expression is the roots-of-unity specialization of the Lagrange
interpolation formula; barycentric forms of Lagrange interpolation are
reviewed in \cite{berrut2004}.  The denominators
\[
x-\omega^i
\]
may be inverted together by the standard batch-inversion technique.  The
implementation therefore replaces \(N\) independent field inversions by
one field inversion and two linear multiplication scans.  The remaining
numerator accumulation is linear in \(N\), and it admits a parallel
reduction.

This representation is common when the surrounding computation already
uses evaluations on a roots-of-unity domain.  For one arbitrary
off-domain point, however, the evaluation requires substantially more
field arithmetic and temporary memory than Horner's rule.  The benchmark
in Section~\ref{sec:results} measures this cost directly.

\subsection{Boolean-Kernel Representation}
\label{sec:boolean-kernel}

We now consider the representation used in this work.  Let
\[
N=2^m
\]
and identify an integer
\[
0\le y<N
\]
with its binary vector
\[
y=(y_0,\ldots,y_{m-1})\in\{0,1\}^m.
\]
For each such vector define
\[
K_y(X)
=
\prod_{i=0}^{m-1}
\left(X^{2^i}+y_i\right).
\]
Every \(K_y\) has degree \(N-1\), and the family
\[
\{K_y : y\in\{0,1\}^m\}
\]
forms a basis of the vector space of polynomials of degree smaller than
\(N\).  We write
\[
P(X)
=
\sum_{y\in\{0,1\}^m}
c_yK_y(X).
\]

The basis differs from the monomial basis in an important point.  Its
elements share a product structure indexed by the bits of \(y\).  At a
fixed point \(x\), the factor associated with bit \(i\) is either
\[
x^{2^i}
\]
or
\[
x^{2^i}+1.
\]
Pairs of coefficients that differ only in this bit can therefore be
combined before the remaining factors are evaluated.

The change of basis from monomial coefficients can be expressed as a
Boolean zeta/Möbius transform and requires \(O(N\log N)\) field additions.
That conversion is not included in the isolated single-point evaluation
benchmark.  As with the Lagrange representation, the benchmark assumes that
the polynomial is already available in the representation under study.

\subsection{Boolean-Kernel Evaluation by Folding}
\label{sec:kernel-fold}

Let
\[
t_i=x^{2^i}.
\]
Consider two terms whose Boolean indices agree except at coordinate \(i\).
After factoring all the common kernel factors, their contribution has the
form
\[
a\,t_i+b(t_i+1).
\]
It can be replaced by the single coefficient
\[
F_{t_i}(a,b)
=
t_i(a+b)+b.
\]
Thus one Boolean coordinate is eliminated by the local transformation
\[
\boxed{
(a,b)\longmapsto t_i(a+b)+b.
}
\]

Starting from the vector
\[
(c_0,c_1,\ldots,c_{N-1}),
\]
the first level combines adjacent pairs using \(t_0=x\).  The next level
combines adjacent results using \(t_1=x^2\), and the process continues with
\[
t_{i+1}=t_i^2
\]
until one field element remains.  The final value is \(P(x)\).

At level \(i\), the number of folds is
\[
\frac{N}{2^{i+1}}.
\]
The total number is therefore
\[
\frac{N}{2}
+\frac{N}{4}
+\cdots+1
=
N-1.
\]
Each fold contains one field multiplication and two field additions.  The
evaluation consequently uses
\[
N-1
\]
multiplications,
\[
2(N-1)
\]
additions, and
\[
m-1
\]
squarings for the sequence \(t_0,\ldots,t_{m-1}\).

Figure~\ref{fig:dependency} compares the dependency graph with the
classical Horner recurrence for a small instance.  The Horner graph is one
chain.  The Boolean-kernel graph has the same logarithmic number of levels
as a balanced reduction tree, with all nodes at one level independent.

\begin{figure}[ht]
\centering
\setlength{\unitlength}{0.82mm}
\begin{picture}(150,58)

% Horner chain
\put(3,52){\makebox(55,4)[c]{\small (a) Horner dependency chain}}
\put(8,40){\circle{6}}
\put(8,40){\makebox(0,0){\scriptsize \(a_7\)}}
\put(22,40){\circle{6}}
\put(22,40){\makebox(0,0){\scriptsize \(r_6\)}}
\put(36,40){\circle{6}}
\put(36,40){\makebox(0,0){\scriptsize \(r_5\)}}
\put(50,40){\circle{6}}
\put(50,40){\makebox(0,0){\scriptsize \(\cdots\)}}
\put(64,40){\circle{6}}
\put(64,40){\makebox(0,0){\scriptsize \(r_0\)}}
\put(11,40){\vector(1,0){8}}
\put(25,40){\vector(1,0){8}}
\put(39,40){\vector(1,0){8}}
\put(53,40){\vector(1,0){8}}

% Kernel tree
\put(88,52){\makebox(56,4)[c]{\small (b) Boolean-kernel reduction}}
\put(82,8){\circle{5}}\put(82,8){\makebox(0,0){\scriptsize \(c_0\)}}
\put(92,8){\circle{5}}\put(92,8){\makebox(0,0){\scriptsize \(c_1\)}}
\put(102,8){\circle{5}}\put(102,8){\makebox(0,0){\scriptsize \(c_2\)}}
\put(112,8){\circle{5}}\put(112,8){\makebox(0,0){\scriptsize \(c_3\)}}
\put(122,8){\circle{5}}\put(122,8){\makebox(0,0){\scriptsize \(c_4\)}}
\put(132,8){\circle{5}}\put(132,8){\makebox(0,0){\scriptsize \(c_5\)}}
\put(142,8){\circle{5}}\put(142,8){\makebox(0,0){\scriptsize \(c_6\)}}
\put(152,8){\circle{5}}\put(152,8){\makebox(0,0){\scriptsize \(c_7\)}}

\put(87,23){\circle{6}}\put(87,23){\makebox(0,0){\scriptsize \(F_0\)}}
\put(107,23){\circle{6}}\put(107,23){\makebox(0,0){\scriptsize \(F_0\)}}
\put(127,23){\circle{6}}\put(127,23){\makebox(0,0){\scriptsize \(F_0\)}}
\put(147,23){\circle{6}}\put(147,23){\makebox(0,0){\scriptsize \(F_0\)}}

\put(97,37){\circle{6}}\put(97,37){\makebox(0,0){\scriptsize \(F_1\)}}
\put(137,37){\circle{6}}\put(137,37){\makebox(0,0){\scriptsize \(F_1\)}}

\put(117,50){\circle{6}}\put(117,50){\makebox(0,0){\scriptsize \(F_2\)}}

\put(83,10){\line(1,3){3.2}}
\put(91,10){\line(-1,3){3.2}}
\put(103,10){\line(1,3){3.2}}
\put(111,10){\line(-1,3){3.2}}
\put(123,10){\line(1,3){3.2}}
\put(131,10){\line(-1,3){3.2}}
\put(143,10){\line(1,3){3.2}}
\put(151,10){\line(-1,3){3.2}}

\put(89,25){\line(2,3){6}}
\put(105,25){\line(-2,3){6}}
\put(129,25){\line(2,3){6}}
\put(145,25){\line(-2,3){6}}

\put(99,39){\line(3,2){15}}
\put(135,39){\line(-3,2){15}}

\end{picture}
\caption{Dependency graphs for \(N=8\).  The classical Horner method
forms a single recurrence.  The Boolean-kernel method has four independent
folds at the first level, two at the second, and one at the last level.
Here \(F_i(a,b)=x^{2^i}(a+b)+b\).}
\label{fig:dependency}
\end{figure}

The two graphs illustrate the principal difference considered in the
performance study.  Horner's method performs fewer additions and has a very
compact implementation, whereas the Boolean-kernel method trades additional
addition work for a much shorter critical path and a large amount of
independent work in the lower levels of the tree.

\section{Goldilocks Field Arithmetic}
\label{sec:goldilocks}

The polynomial evaluation methods of Section~\ref{sec:poly-eval} ultimately
depend on the cost of field multiplication.  We therefore separate the
choice of polynomial representation from the choice of modular
multiplication algorithm.

The experiments use the Goldilocks prime
\[
p=2^{64}-2^{32}+1.
\]
A field element is represented by one machine word.  Two multiplication
strategies are considered.  The first exploits the special form of the
modulus.  The second uses Montgomery representation with radix
\[
R=2^{64}.
\]
Both strategies are retained in the experiments so that the effect of the
modular multiplication algorithm can be measured independently of the
polynomial representation.

\subsection{Reduction Specialized to the Goldilocks Modulus}
\label{sec:goldilocks-special}

Let
\[
z=ab
\]
be the \(128\)-bit product of two canonical residues
\[
0\le a,b<p.
\]
Write
\[
z=z_{\mathrm{lo}}+2^{64}z_{\mathrm{hi}},
\]
where
\[
0\le z_{\mathrm{lo}},z_{\mathrm{hi}}<2^{64}.
\]

For the Goldilocks prime,
\[
2^{64}\equiv 2^{32}-1\pmod p.
\]
Define
\[
\varepsilon=2^{32}-1.
\]
Then
\[
z
\equiv
z_{\mathrm{lo}}+\varepsilon z_{\mathrm{hi}}
\pmod p.
\]

A direct application of this congruence already reduces the \(128\)-bit
product to a value close to one machine word.  A more useful form is
obtained by splitting the high half as
\[
z_{\mathrm{hi}}
=
h_0+2^{32}h_1,
\qquad
0\le h_0,h_1<2^{32}.
\]
Since
\[
2^{32}\varepsilon
=
2^{64}-2^{32}
\equiv -1
\pmod p,
\]
we obtain
\[
\varepsilon z_{\mathrm{hi}}
\equiv
h_0\varepsilon-h_1
\pmod p.
\]
Therefore
\[
z
\equiv
z_{\mathrm{lo}}
-h_1
+h_0(2^{32}-1)
\pmod p.
\]

This expression involves only the \(128\)-bit product \(ab\), extraction of
the two \(32\)-bit halves of \(z_{\mathrm{hi}}\), shifts, additions, and
subtractions, followed by a small number of corrections to return to the
canonical interval.

The implementation used in the benchmark follows this reduction rather than
a generic integer remainder instruction.

\paragraph{Algorithm 1. Goldilocks multiplication.}

\[
\begin{array}{ll}
\textbf{Input:} & a,b\in[0,p),\\
\textbf{Output:} & ab\bmod p.
\end{array}
\]

\[
\begin{array}{rcl}
z &\leftarrow& a\cdot b \qquad\text{(full \(128\)-bit product)},\\
z_{\mathrm{lo}} &\leftarrow& z\bmod 2^{64},\\
z_{\mathrm{hi}} &\leftarrow& \lfloor z/2^{64}\rfloor,\\
h_0 &\leftarrow& z_{\mathrm{hi}}\bmod 2^{32},\\
h_1 &\leftarrow& \lfloor z_{\mathrm{hi}}/2^{32}\rfloor,\\
u &\leftarrow& z_{\mathrm{lo}}-h_1+h_0(2^{32}-1),\\
u &\leftarrow& \text{canonical reduction of }u\text{ modulo }p,\\
\textbf{return}&&u.
\end{array}
\]

The exact instruction sequence generated from this formula depends on the
compiler and target architecture.  For this reason the multiplication is
measured separately before its effect on the complete polynomial evaluator
is considered.

\subsection{Montgomery Representation}
\label{sec:montgomery}

Montgomery multiplication replaces division by \(p\) with arithmetic modulo
a power of two.  Let
\[
R=2^{64}.
\]
Because \(p\) is odd,
\[
\gcd(p,R)=1,
\]
and \(R\) is invertible modulo \(p\).

A canonical field element \(a\) is represented in Montgomery form by
\[
\widetilde{a}=aR\bmod p.
\]
Multiplication of two Montgomery residues produces
\[
\widetilde{a}\widetilde{b}
=
abR^2
\]
and is followed by a Montgomery reduction, which removes one factor of
\(R\):
\[
\operatorname{REDC}
\left(
\widetilde{a}\widetilde{b}
\right)
=
abR\bmod p.
\]

For the modulus considered here, the constants are
\[
R\bmod p
=
2^{32}-1
=
\mathtt{0x00000000ffffffff},
\]
\[
n'=-p^{-1}\bmod R
=
\mathtt{0xfffffffeffffffff},
\]
and
\[
R^2\bmod p
=
\mathtt{0xfffffffe00000001}.
\]

The value \(R^2\bmod p\) is used to convert a canonical residue into
Montgomery form with one Montgomery multiplication.

\subsection{Montgomery Reduction}
\label{sec:montgomery-redc}

Let
\[
0\le T<pR.
\]
The Montgomery reduction computes
\[
TR^{-1}\bmod p
\]
without an integer division by \(p\).

Using
\[
n'=-p^{-1}\bmod R,
\]
define
\[
m=(T\bmod R)n'\bmod R.
\]
By construction,
\[
T+mp
\equiv 0\pmod R,
\]
so the quotient
\[
u=\frac{T+mp}{R}
\]
is an integer.  The result satisfies
\[
u\equiv TR^{-1}\pmod p.
\]
A final conditional subtraction returns the canonical representative.

\paragraph{Algorithm 2. Montgomery reduction.}

\[
\begin{array}{ll}
\textbf{Input:} & T,\quad 0\le T<pR,\\
\textbf{Output:} & TR^{-1}\bmod p.
\end{array}
\]

\[
\begin{array}{rcl}
m &\leftarrow& (T\bmod R)n'\bmod R,\\
u &\leftarrow& (T+mp)/R,\\
\textbf{if }u\ge p && u\leftarrow u-p,\\
\textbf{return}&&u.
\end{array}
\]

For a \(64\)-bit radix, the sum
\[
T+mp
\]
may require \(129\) bits.  The implementation therefore records the carry
from the \(128\)-bit addition explicitly before extracting the high word.

\subsection{Montgomery Multiplication}
\label{sec:montgomery-mul}

If \(\widetilde{a}\) and \(\widetilde{b}\) are already stored in Montgomery
form, their product is computed as
\[
\operatorname{MontMul}
(\widetilde{a},\widetilde{b})
=
\operatorname{REDC}
(\widetilde{a}\widetilde{b}).
\]

\paragraph{Algorithm 3. Montgomery multiplication.}

\[
\begin{array}{ll}
\textbf{Input:} & \widetilde{a}=aR\bmod p,\quad
                  \widetilde{b}=bR\bmod p,\\
\textbf{Output:} & abR\bmod p.
\end{array}
\]

\[
\begin{array}{rcl}
T &\leftarrow& \widetilde{a}\widetilde{b},\\
u &\leftarrow& \operatorname{REDC}(T),\\
\textbf{return}&&u.
\end{array}
\]

Conversion from canonical to Montgomery form is performed as
\[
\widetilde{a}
=
\operatorname{REDC}
\left(
a(R^2\bmod p)
\right),
\]
while conversion back to the canonical representation is
\[
a
=
\operatorname{REDC}(\widetilde{a}).
\]

In the multiplication benchmark, operands are kept in their native
representation throughout the timed region.  Thus the specialized
Goldilocks implementation receives canonical residues, while the
Montgomery implementation receives Montgomery residues.  Conversion is
measured separately when it is relevant to an application and is not
charged once per multiplication.

\subsection{Kernel Fold Under the Two Arithmetic Models}
\label{sec:fold-arithmetic}

The Boolean-kernel reduction uses
\[
F_t(a,b)=t(a+b)+b.
\]
With canonical Goldilocks residues, the implementation can combine the
final addition by \(b\) with the reduction of the \(128\)-bit product.
This gives a specialized field operation rather than a literal sequence of
one modular addition, one modular multiplication, and a second modular
addition.

In Montgomery form the same mathematical operation is
\[
\widetilde{F}_t(\widetilde{a},\widetilde{b})
=
\operatorname{MontMul}
\left(
\widetilde{t},
\widetilde{a}+\widetilde{b}
\right)
+
\widetilde{b}.
\]
All operands remain in Montgomery form, so no conversion is needed between
levels of the reduction tree.

Table~\ref{tab:field-mul-structure} summarizes the two multiplication
strategies before measurements.

\begin{table}[ht]
\centering
\caption{Structure of the two Goldilocks multiplication strategies.}
\label{tab:field-mul-structure}
\begin{tabular}{lll}
\toprule
Method & Representation & Main reduction operation\\
\midrule
Goldilocks special &
canonical residue &
use \(2^{64}\equiv2^{32}-1\pmod p\)\\

Montgomery &
\(aR\bmod p\) &
\((T+mp)/2^{64}\) with \(m=T_0n'\)\\
\bottomrule
\end{tabular}
\end{table}

At this point there is no architectural conclusion to draw from the
operation formulas alone.  The specialized reduction uses the algebraic
form of the Goldilocks modulus; Montgomery reduction introduces additional
word multiplications but has a regular reduction structure.  The generated
instruction sequences and the complete evaluation times are therefore
examined experimentally in the following sections.

\section{Processor and Measurement Methods}
\label{sec:methods}

\subsection{Latency and Throughput Definitions}
\label{sec:latency-throughput}

The execution time of a polynomial evaluator depends on two distinct
properties of the instruction stream.  The first is the delay introduced by
a dependency chain.  The second is the rate at which independent
instructions of the same class can be issued.

A \emph{dependency chain} is a sequence of operations in which each result
is required by the next operation.  The classical Horner recurrence is the
simplest example in this study:
\[
r_i=a_i+xr_{i+1}.
\]
Even if several multiplication units are available, the next multiplication
cannot begin before the preceding value \(r_{i+1}\) has been produced.

The \emph{latency} of an operation is the delay that it contributes to such
a dependency chain.  The \emph{throughput} is the rate at which independent
operations can be completed when enough independent operands are available.
For the algorithms considered here, neither quantity alone is sufficient.
A method with a longer local operation sequence can still execute faster if
many instances of that sequence are independent.

The Boolean-kernel reduction makes this distinction explicit.  At level
\(i\), all
\[
\frac{N}{2^{i+1}}
\]
folds are independent, while the transition from one level to the next is a
true dependency.  The critical path therefore grows with the number of
levels, \(m=\log_2 N\), rather than with the number of coefficients.

\subsection{Architecture Details}
\label{sec:architecture}

The measurements reported in this study were obtained on an Apple M1
processor.  The comparison follows the same general processor-oriented
principle used in earlier work on polynomial evaluation: arithmetic count
alone is insufficient when evaluation schemes expose different dependency
graphs \cite{ewart2020}.  The M1 is an out-of-order superscalar processor,
so several integer operations can be in flight simultaneously when the
instruction scheduler finds independent work.  Horner evaluation, blocked
Horner evaluation, and the Boolean-kernel tree expose very different
amounts of instruction-level and task-level parallelism.

We do not attempt to model the complete microarchitecture.  Only the
features required to interpret the measurements are considered: independent
integer multiplication, dependency depth, memory traffic, allocation, and
the scheduling of independent subtrees across CPU workers.

The field implementation also matters at the instruction level.  A
Goldilocks multiplication specialized to
\[
p=2^{64}-2^{32}+1
\]
uses one full-width product followed by shifts, additions, and
subtractions.  A Montgomery multiplication uses the initial product and an
additional product involving the reduction coefficient.  Figure
\ref{fig:field-dag} shows the corresponding dependency structure
schematically.

\begin{figure}[ht]
\centering
\setlength{\unitlength}{0.83mm}
\begin{picture}(150,65)

\put(5,59){\makebox(60,4)[c]{\small (a) Goldilocks special reduction}}
\put(95,59){\makebox(48,4)[c]{\small (b) Montgomery reduction}}

% Goldilocks DAG
\put(12,48){\framebox(18,7){\(a\times b\)}}
\put(12,35){\framebox(18,7){split hi/lo}}
\put(7,21){\framebox(28,7){split \(h_0,h_1\)}}
\put(7,8){\framebox(28,7){shift/add/sub}}
\put(12,48){\vector(0,-1){6}}
\put(21,35){\vector(0,-1){7}}
\put(21,21){\vector(0,-1){6}}

\put(43,48){\framebox(18,7){\(a+b\)}}
\put(43,35){\framebox(18,7){\(t(a+b)\)}}
\put(43,21){\framebox(18,7){reduce}}
\put(43,8){\framebox(18,7){\(+b\)}}
\put(52,48){\vector(0,-1){6}}
\put(52,35){\vector(0,-1){7}}
\put(52,21){\vector(0,-1){6}}

% Montgomery DAG
\put(94,48){\framebox(20,7){\(T=a b\)}}
\put(94,35){\framebox(20,7){\(m=T_0n'\)}}
\put(92,21){\framebox(24,7){\(T+mp\)}}
\put(94,8){\framebox(20,7){high word}}
\put(104,48){\vector(0,-1){6}}
\put(104,35){\vector(0,-1){7}}
\put(104,21){\vector(0,-1){6}}

\put(126,48){\framebox(20,7){\(\tilde a+\tilde b\)}}
\put(126,35){\framebox(20,7){MontMul}}
\put(126,21){\framebox(20,7){\(+\tilde b\)}}
\put(136,48){\vector(0,-1){6}}
\put(136,35){\vector(0,-1){7}}

\end{picture}
\caption{Schematic dependency graphs for the two field-arithmetic paths.
The left side shows the specialized Goldilocks product reduction and its
use in the kernel fold; the right side shows the Montgomery reduction path.
The diagram represents dependencies rather than exact machine
instructions.}
\label{fig:field-dag}
\end{figure}

\subsection{Exact Arithmetic}
\label{sec:exact-arithmetic}

Unlike floating-point polynomial evaluation, the present computation does
not involve a numerical approximation error.  Every benchmarked method is
required to return exactly the same element of \(\mathbb{F}_p\).

Correctness is checked before performance measurements.  The test suite
compares the optimized Goldilocks multiplication with a reference modular
multiplication, compares the fused Boolean-kernel fold with its direct field
expression, verifies the monomial-to-kernel change of basis by evaluating
the same polynomial in both representations, and verifies the Lagrange
implementation against the monomial reference.

The Montgomery implementation is subjected to the same requirement.
Canonical residues are converted to Montgomery form, multiplied there, and
converted back before comparison with the reference result.  This prevents
a performance result from being accepted merely because two implementations
use different internal encodings.

\subsection{Measurement Tools}
\label{sec:measurement-tools}

The benchmarks are implemented with Criterion.  The principal measurements
use a warm-up period before sampling and report an interval together with a
central estimate.  The benchmark harness applies \texttt{black\_box} to
inputs and results in order to prevent the compiler from replacing the
measured computation with a constant result.

Microbenchmarks are used for the field operations and the kernel fold.
Complete polynomial evaluators are measured separately.  This separation is
important: an improvement in the latency of one multiplication need not
translate linearly to the full evaluator once memory traffic, worker
scheduling, and allocation are included.

For cross-method comparisons, all implementations are compiled from the same
repository revision and executed on the same machine.  A ratio is formed
only from measurements collected under the same field backend and benchmark
configuration.

\subsection{Programming Model}
\label{sec:programming-model}

The implementation is written in Rust.  Parallel work is expressed with
Rayon.  The monomial baseline divides the coefficient vector into blocks,
evaluates those blocks independently, and combines the resulting values.
The Boolean-kernel implementation divides the reduction tree into
power-of-two subtrees.  Each worker reduces one complete subtree locally,
thereby avoiding a global synchronization after every tree level.

The optimized kernel evaluator additionally reuses worker-local scratch
storage.  The first tree level is written directly from the immutable input
slice into a half-size scratch buffer, after which the remaining local
levels are performed in place.

This implementation preserves the arithmetic graph derived in
Section~\ref{sec:kernel-fold}; it changes only the scheduling and memory
traffic of that graph.

\subsection{Processor and Compiler}
\label{sec:processor-compiler}

Table~\ref{tab:platform} records the platform used for the current
measurement campaign.  The values are generated from the machine on which
this manuscript is built.

\begin{table}[ht]
\centering
\caption{Processor, compiler, and benchmark environment.}
\label{tab:platform}
\begin{tabular}{ll}
\toprule
Item & Configuration\\
\midrule
Processor & Apple M1\\
Physical CPU cores & 8\\
Logical CPU cores & 8\\
Memory & 16 GiB\\
Rust compiler & rustc 1.96.0-nightly (55e86c996 2026-04-02)\\
Cargo & cargo 1.96.0-nightly (888f67534 2026-03-30)\\
Operating system & Darwin 25.6.0 arm64\\
Parallel runtime & Rayon\\
\bottomrule
\end{tabular}
\end{table}

The benchmarks are compiled in release mode.  The repository configuration
uses link-time optimization for the performance runs.  No result in the
paper is taken from a debug build.

\subsection{Validation}
\label{sec:validation}

A simple consistency check is available before examining the final results.
The classical Horner method forms one dependency chain, whereas the
Boolean-kernel evaluator exposes independent folds at every level.  The
latter should therefore benefit more strongly from multiple workers once
the problem size is large enough to amortize scheduling overhead.

At small sizes the opposite effect is possible: the overhead of creating
parallel tasks and managing scratch storage can be comparable to the useful
work.  The expected behavior is consequently a crossover rather than a
uniform speedup at all sizes.

The field backend provides a second validation point.  If the
specialized Goldilocks reduction has a lower isolated multiplication cost
than Montgomery reduction, this advantage should be visible most clearly in
the serial Horner benchmark, whose execution is dominated by a chain of
multiplications.  In the Boolean-kernel evaluator the effect can be reduced
or amplified by memory traffic and parallel scheduling.  The measurements
in the following section are interpreted against these two expectations.


\section{Results}
\label{sec:results}

Table~\ref{tab:results-summary} gives a compact view of the principal
measurements.  Two experiments are reported separately.  The first is the
cross-representation sweep, in which all polynomial representations are
measured in one benchmark campaign with the same Goldilocks arithmetic
backend.  The second is the backend experiment, in which specialized
Goldilocks arithmetic and Montgomery arithmetic are compared for the same
evaluation algorithm.  Absolute timings are not compared across these two
campaigns.

\begin{table}[ht]
\centering
\caption{Summary of the main measurements at \(N=2^{20}\).  Cross-basis
times come from one common benchmark sweep.  Backend ratios are the median
of five independent same-run comparisons.}
\label{tab:results-summary}
\begin{tabular}{lll}
\toprule
Experiment & Quantity & Result\\
\midrule
Cross-basis &
Boolean-kernel &
1.0560 ms\\
Cross-basis &
parallel monomial &
1.4116 ms\\
Cross-basis &
kernel / parallel monomial &
1.337\(\times\)\\
Backend &
Montgomery / special, kernel &
1.2777\(\times\)\\
Backend &
Montgomery / special, monomial &
1.0591\(\times\)\\
\bottomrule
\end{tabular}
\end{table}

The first result is that the ranking of field multipliers and the ranking
of complete evaluators are not the same question.  In the isolated
multiplication benchmark, specialized Goldilocks multiplication and
Montgomery multiplication have nearly identical central estimates:
\(1.9107\) ns and \(1.9162\) ns, respectively.  Their measurement intervals
overlap.  No meaningful advantage can therefore be assigned to either
method from the isolated multiplication benchmark alone.

The kernel fold behaves differently.  The specialized operation
\[
t(a+b)+b
\]
takes \(3.0292\) ns in the microbenchmark, compared with \(3.7952\) ns for
the Montgomery implementation.  The ratio is \(1.2529\times\).  The
difference appears only after the multiplication has been embedded in the
compound fold operation, where the form of the Goldilocks modulus can be
used together with the surrounding additions.

\subsection{Field Backend and Complete Evaluators}
\label{sec:backend-results}

To test whether the microbenchmark difference survives in the complete
evaluator, the four full \(N=2^{20}\) configurations were repeated in
five independent Criterion invocations.  Table~\ref{tab:backend-stability}
reports the central estimate from each invocation.

\begin{table}[ht]
\centering
\caption{Repeated full-evaluator measurements at \(N=2^{20}\), in
microseconds.  M/S denotes Montgomery time divided by specialized
Goldilocks time within the same invocation.}
\label{tab:backend-stability}
\begin{tabular}{rrrrrr}
\toprule
Run & \(K_S\) & \(K_M\) & \(M_S\) & \(M_M\) & Kernel M/S\\
\midrule
1 & 883.61 & 1172.40 & 1453.50 & 1531.00 & 1.327\\
2 & 989.03 & 1227.70 & 1559.40 & 1654.20 & 1.241\\
3 & 1080.90 & 1159.90 & 1537.00 & 1901.10 & 1.073\\
4 & 1036.00 & 1323.70 & 2208.30 & 1666.00 & 1.278\\
5 & 952.15 & 1277.60 & 1607.80 & 1702.80 & 1.342\\
\midrule
Median &
989.03 &
1227.70 &
1559.40 &
1666.00 &
1.278\\
\bottomrule
\end{tabular}
\end{table}

The kernel result is stable in sign across all five invocations.  The
Montgomery implementation is slower in every run, with a median same-run
ratio of
\[
1.2777\times.
\]
The absolute specialized-kernel estimates vary from
\(883.61\,\mu\mathrm{s}\) to
\(1080.90\,\mu\mathrm{s}\), which is why the comparison is made
within each run rather than by combining absolute times collected at
different moments.

The monomial backend comparison is less stable.  The same-run ratios range
from
\[
0.7544\times
\quad\text{to}\quad
1.2369\times,
\]
and one invocation reverses the ordering.  The data therefore do not
support a claim that either arithmetic backend is systematically faster for
the complete blocked monomial evaluator.  This contrasts with the kernel
case and suggests that the specialized advantage is tied to the fold
structure rather than to scalar field multiplication alone.

\subsection{Cross-Representation Evaluation}
\label{sec:cross-representation}

We next compare the polynomial representations using one common Goldilocks
backend.  At \(N=2^{20}\), the optimized Boolean-kernel evaluator requires
\[
1.0560\ \mathrm{ms},
\]
while the parallel monomial implementation requires
\[
1.4116\ \mathrm{ms}.
\]
The ratio is
\[
\frac{1.4116}{1.0560}
=
1.337.
\]
Equivalently, the Boolean-kernel evaluator uses
\[
25.2\%
\]
less wall time for this isolated evaluation primitive.

The serial Horner method takes
\[
9.0402\ \mathrm{ms},
\]
which is
\[
8.561\times
\]
the kernel time.  Evaluation from the roots-of-unity Lagrange
representation is more expensive for an arbitrary off-domain point:
the parallel implementation takes
\[
17.1990\ \mathrm{ms},
\]
and the serial implementation
\[
21.3890\ \mathrm{ms}.
\]
These numbers characterize only arbitrary single-point evaluation from the
given representation; they do not imply that evaluation-form
representations are generally inefficient in workloads that operate
directly on the evaluation domain.

\subsection{Scaling with the Polynomial Size}
\label{sec:scaling}

Figure~\ref{fig:scaling} reports the complete cross-basis sweep from
\(N=2^{10}\) to \(N=2^{20}\).  The vertical axis is logarithmic in time.
The figure plays the same role as the global result plots in the reference
performance study: it shows the complete set of principal algorithms before
the individual measurements are interpreted.

\begin{figure}[ht]
\centering
\setlength{\unitlength}{1mm}
\begin{picture}(154,72)
\put(16,8){\vector(1,0){102}}
\put(16,8){\vector(0,1){58}}
\put(2,64){\makebox(14,4)[r]{\scriptsize time}}
\put(116,2){\makebox(12,4)[l]{\scriptsize \(N\)}}
\put(20,7){\line(0,1){2}}
\put(16,1){\makebox(8,4)[c]{\scriptsize \(2^{10}\)}}
\put(38,7){\line(0,1){2}}
\put(34,1){\makebox(8,4)[c]{\scriptsize \(2^{12}\)}}
\put(56,7){\line(0,1){2}}
\put(52,1){\makebox(8,4)[c]{\scriptsize \(2^{14}\)}}
\put(74,7){\line(0,1){2}}
\put(70,1){\makebox(8,4)[c]{\scriptsize \(2^{16}\)}}
\put(92,7){\line(0,1){2}}
\put(88,1){\makebox(8,4)[c]{\scriptsize \(2^{18}\)}}
\put(110,7){\line(0,1){2}}
\put(106,1){\makebox(8,4)[c]{\scriptsize \(2^{20}\)}}
\put(15,15.1){\line(1,0){2}}
\put(1,13.6){\makebox(13,3)[r]{\scriptsize 10 us}}
\put(15,29.2){\line(1,0){2}}
\put(1,27.7){\makebox(13,3)[r]{\scriptsize 100 us}}
\put(15,43.3){\line(1,0){2}}
\put(1,41.8){\makebox(13,3)[r]{\scriptsize 1 ms}}
\put(15,57.3){\line(1,0){2}}
\put(1,55.8){\makebox(13,3)[r]{\scriptsize 10 ms}}
\put(20,10.0){\circle*{1.8}}
\put(38,18.3){\circle*{1.8}}
\put(56,26.0){\circle*{1.8}}
\put(74,28.6){\circle*{1.8}}
\put(92,35.1){\circle*{1.8}}
\put(110,43.6){\circle*{1.8}}
\qbezier(20,10.0)(29.0,14.15)(38,18.3)
\qbezier(38,18.3)(47.0,22.15)(56,26.0)
\qbezier(56,26.0)(65.0,27.3)(74,28.6)
\qbezier(74,28.6)(83.0,31.85)(92,35.1)
\qbezier(92,35.1)(101.0,39.35)(110,43.6)
\put(20,14.2){\circle{2.2}}
\put(38,22.7){\circle{2.2}}
\put(56,26.9){\circle{2.2}}
\put(74,31.2){\circle{2.2}}
\put(92,37.4){\circle{2.2}}
\put(110,45.4){\circle{2.2}}
\qbezier(20,14.2)(29.0,18.45)(38,22.7)
\qbezier(38,22.7)(47.0,24.799999999999997)(56,26.9)
\qbezier(56,26.9)(65.0,29.049999999999997)(74,31.2)
\qbezier(74,31.2)(83.0,34.3)(92,37.4)
\qbezier(92,37.4)(101.0,41.4)(110,45.4)
\put(18.9,13.200000000000001){\framebox(2.2,2.2){}}
\put(36.9,21.7){\framebox(2.2,2.2){}}
\put(54.9,30.4){\framebox(2.2,2.2){}}
\put(72.9,38.6){\framebox(2.2,2.2){}}
\put(90.9,47.1){\framebox(2.2,2.2){}}
\put(108.9,55.6){\framebox(2.2,2.2){}}
\qbezier(20,14.3)(29.0,18.55)(38,22.8)
\qbezier(38,22.8)(47.0,27.15)(56,31.5)
\qbezier(56,31.5)(65.0,35.6)(74,39.7)
\qbezier(74,39.7)(83.0,43.95)(92,48.2)
\qbezier(92,48.2)(101.0,52.45)(110,56.7)
\put(20,29.8){\circle*{1.8}}
\put(38,34.8){\circle*{1.8}}
\put(56,38.2){\circle*{1.8}}
\put(74,44.6){\circle*{1.8}}
\put(92,52.5){\circle*{1.8}}
\put(110,60.7){\circle*{1.8}}
\qbezier(20,29.8)(29.0,32.3)(38,34.8)
\qbezier(38,34.8)(47.0,36.5)(56,38.2)
\qbezier(56,38.2)(65.0,41.400000000000006)(74,44.6)
\qbezier(74,44.6)(83.0,48.55)(92,52.5)
\qbezier(92,52.5)(101.0,56.6)(110,60.7)
\put(20,19.8){\circle{2.2}}
\put(38,28.1){\circle{2.2}}
\put(56,36.6){\circle{2.2}}
\put(74,45.1){\circle{2.2}}
\put(92,53.5){\circle{2.2}}
\put(110,62.0){\circle{2.2}}
\qbezier(20,19.8)(29.0,23.950000000000003)(38,28.1)
\qbezier(38,28.1)(47.0,32.35)(56,36.6)
\qbezier(56,36.6)(65.0,40.85)(74,45.1)
\qbezier(74,45.1)(83.0,49.3)(92,53.5)
\qbezier(92,53.5)(101.0,57.75)(110,62.0)
\put(122,60){\circle*{1.8}}
\put(126,58.5){\makebox(24,3)[l]{\scriptsize \(K\)}}
\put(122,53){\circle{2.2}}
\put(126,51.5){\makebox(24,3)[l]{\scriptsize \(M_p\)}}
\put(120.9,44.9){\framebox(2.2,2.2){}}
\put(126,44.5){\makebox(24,3)[l]{\scriptsize \(M_h\)}}
\put(122,39){\circle*{1.8}}
\put(126,37.5){\makebox(24,3)[l]{\scriptsize \(L_p\)}}
\put(122,32){\circle{2.2}}
\put(126,30.5){\makebox(24,3)[l]{\scriptsize \(L_s\)}}
\end{picture}
\caption{Single-point evaluation time as a function of \(N\) for the five
implementations in the common cross-basis benchmark.  \(K\): optimized
Boolean-kernel; \(M_p\): parallel blocked monomial; \(M_h\): serial Horner;
\(L_p\): parallel Lagrange; \(L_s\): serial Lagrange.  The vertical
placement is logarithmic in measured time.}
\label{fig:scaling}
\end{figure}

The small-size regime is affected strongly by parallel scheduling and fixed
overheads.  At \(N=2^{10}\), the kernel evaluator takes approximately half
the time of the parallel monomial evaluator, but the absolute difference is
only a few microseconds.  At \(N=2^{14}\), the two parallel methods become
closer, with a ratio of approximately
\[
1.169\times.
\]
For larger inputs the separation grows again.  At \(N=2^{18}\), the
parallel monomial evaluator is
\[
1.447\times
\]
the kernel time, and at \(N=2^{20}\) the ratio is
\[
1.337\times.
\]

The serial Horner curve grows linearly with a long dependency chain and
separates rapidly from the two parallel implementations.  The two Lagrange
curves have the largest constants because arbitrary off-domain evaluation
includes denominator construction, batch inversion, and a full linear
accumulation.

\subsection{Interpretation}
\label{sec:results-interpretation}

The measurements support two separate conclusions.  First, the
Boolean-kernel representation gives a lower single-point evaluation time
than the fair parallel monomial baseline in the common \(N=2^{20}\)
benchmark.  Second, this advantage cannot be explained by a generally
faster scalar field multiplication: isolated Goldilocks and Montgomery
multiplication are nearly tied, whereas the specialized Boolean-kernel fold
and the complete specialized kernel evaluator show a clear advantage.

The result is therefore an interaction between representation and
implementation.  The Boolean-kernel basis exposes a balanced reduction
tree, and the Goldilocks modulus permits the arithmetic inside each fold to
be specialized as one compound operation.  The benchmark does not establish
a corresponding whole-prover speedup.  Whether the gain survives in a
proof system depends on how long the polynomial can remain in the
Boolean-kernel representation and on the cost of any required basis
conversion.

\section{Conclusions}
\label{sec:conclusions}

This article compared single-point evaluation of univariate polynomials in
three representations on an Apple M1 processor.  The Boolean-kernel basis
leads to a balanced reduction tree in which all folds at one level are
independent.  In the common cross-representation benchmark at
\(N=2^{20}\), this evaluator required \(1.0560\) ms, compared with
\(1.4116\) ms for the blocked parallel monomial implementation.  The
difference corresponds to \(25.2\%\) less wall time for the isolated
evaluation operation.

The arithmetic-backend experiment helps identify the source of this
difference.  Specialized Goldilocks multiplication and Montgomery
multiplication have almost identical isolated multiplication times.  The
fold operation does not: the specialized Goldilocks fold is faster than
its Montgomery counterpart, and the same ordering is observed in all five
independent runs of the complete kernel evaluator.  No similarly stable
ordering is observed for the blocked monomial evaluator.  The performance
difference is therefore associated with the combination of the
Boolean-kernel reduction structure and the arithmetic used inside the
fold, rather than with a universal advantage of one scalar multiplication
routine.

These measurements concern a single processor and an isolated polynomial
operation.  They do not establish an end-to-end speedup for a proof system.
Such a speedup depends on the representation used by the surrounding
protocol, the frequency of basis conversion, and the fraction of prover
time spent in the evaluated primitive.  A monomial-to-kernel conversion
requires \(O(N\log N)\) field additions, so repeated conversion can erase
the benefit measured here.

Further work should therefore address two questions.  The first is
architectural: the same evaluators should be measured on additional
out-of-order CPUs and, where appropriate, on SIMD and accelerator
implementations.  The second is algorithmic: the Boolean-kernel
representation should be integrated into a complete polynomial-commitment
or proof-system pipeline in which conversion, commitment, and opening costs
are measured together.  These experiments would determine whether the
single-point advantage observed here persists beyond the isolated
evaluation kernel.


\bibliographystyle{plain}
\bibliography{references}

\end{document}
